% ============================================================
%  Chapitre 4 — Analyse de la version de référence (mono-cœur)
% ============================================================
\chapter{Version de référence : analyse mono-cœur}
\label{chap:baseline}

\section{Instrumentation du code original (AoS)}

\subsection{Démarche}

La première étape consiste à \emph{instrumenter} le code de référence (\texttt{00_src/})
sans modifier sa logique, afin d'établir une ligne de base de performance.
On encadre chaque phase de la boucle principale avec des chronomètres
\texttt{std::chrono::steady\_clock} et on accumule les durées sur 100 itérations
avant de les afficher sur \texttt{stdout}.

La structure ajoutée est la suivante :

\begin{lstlisting}[language=C++, caption={Structure d'accumulation des temps (ant\_simu.cpp)}]
struct TimingStats {
    double t_ants = 0., t_evap = 0., t_update = 0., t_render = 0.;
    int    count  = 0;
    void print() const {
        double inv = 1e3 / count;
        std::cout << "Fourmis : "    << t_ants  * inv << " ms/it\n"
                  << "Evaporation : "<< t_evap  * inv << " ms/it\n"
                  << "Update : "     << t_update* inv << " ms/it\n"
                  << "Rendu SDL : "  << t_render* inv << " ms/it\n";
    }
};
\end{lstlisting}

\subsection{Résultats — version AoS séquentielle}

Les mesures ci-dessous proviennent de runs instrumentés de
\texttt{src\_v1\_timing/} exportés dans \texttt{profiling.csv}, puis agrégés par
\texttt{profile\_stats.exe}. Deux campagnes sont comparées :
\begin{itemize}
    \item un \textbf{run long} de \textbf{14184 itérations} ;
    \item un \textbf{run court} de \textbf{3851 itérations}.
\end{itemize}

On relève également deux indicateurs métier utiles :
\begin{itemize}
    \item première nourriture ramenée au nid à l'itération \textbf{1747} dans les deux runs ;
    \item quantité finale de nourriture : \textbf{46062} (run long) vs \textbf{1221} (run court).
\end{itemize}

\begin{table}[h]
\centering
\caption{Profiling par phase — version AoS séquentielle (1 cœur, run long)}
\label{tab:aos_timing}
\begin{tabular}{lrr}
\hline
\textbf{Phase}      & \textbf{Temps (ms/it)} & \textbf{Part (\%)} \\
\hline
Avancement fourmis  & 4.1713   & 4.05  \\
Évaporation         & 0.6541   & 0.64  \\
Update (swap)       & 0.0119   & 0.01  \\
Rendu SDL           & 98.1038  & 95.30 \\
\hline
\textbf{TOTAL}      & \textbf{102.9411} & 100.0 \\
\hline
\end{tabular}
\end{table}

\begin{table}[h]
\centering
\caption{Comparaison inter-runs (14184 vs 3851 itérations)}
\label{tab:aos_timing_compare}
\begin{tabular}{lrrrr}
\hline
\textbf{Phase} & \textbf{Run long (ms/it)} & \textbf{Run court (ms/it)} & \textbf{Delta abs.} & \textbf{Delta (\%)} \\
\hline
Avancement fourmis & 4.1713 & 5.0230 & +0.8517 & +20.4 \\
Évaporation        & 0.6541 & 0.7179 & +0.0638 & +9.8  \\
Update             & 0.0119 & 0.0134 & +0.0015 & +12.6 \\
Rendu SDL          & 98.1038 & 92.0667 & -6.0371 & -6.2 \\
\hline
\textbf{TOTAL}     & \textbf{102.9411} & \textbf{97.8209} & \textbf{-5.1202} & \textbf{-5.0} \\
\hline
\end{tabular}
\end{table}

\subsection{Analyse}

Le goulot dominant reste \textbf{le rendu SDL} dans les deux campagnes
(\(94.12\%\) à \(95.30\%\) du temps total). La conclusion principale est donc
\textbf{stable} : les mesures end-to-end sont majoritairement pilotées par la
chaîne graphique et par le scheduling système.

La comparaison inter-runs montre toutefois une variabilité non négligeable :
\begin{itemize}
    \item le temps total moyen diminue de \(\approx 5\%\) sur le run court ;
    \item à l'inverse, le coût \emph{algorithmique} (hors rendu) augmente.
\end{itemize}

En isolant le coût algorithmique (sans rendu), on obtient :
\[
T_{\mathrm{calc}} = T_{\mathrm{ants}} + T_{\mathrm{evap}} + T_{\mathrm{update}}
\approx 4.1713 + 0.6541 + 0.0119 = 4.8373\ \mathrm{ms/it}
\]
pour le run long, et
\[
T_{\mathrm{calc,\ court}} \approx 5.0230 + 0.7179 + 0.0134 = 5.7543\ \mathrm{ms/it}.
\]
soit une cadence théorique d'environ
\[
\frac{1}{0.0048373} \approx 206\ \mathrm{it/s}
\]
sur le run long (contre \(\approx 9.7\ \mathrm{it/s}\) en mode affiché), et
\(\approx 174\ \mathrm{it/s}\) sur le run court.

Cette différence confirme que les conclusions sur les phases chronophages doivent
être tirées \textbf{en priorité sur les parts relatives} et sur des campagnes
répétées, plutôt que sur une seule valeur absolue de latence.

Conclusion méthodologique : pour évaluer la scalabilité OpenMP/MPI, il faut
utiliser des mesures \textbf{sans rendu} (ou avec rendu décimé), sinon les
accélérations observées reflètent surtout le sous-système graphique plutôt que
l'algorithme ACO.

\section{Vectorisation par passage SoA}

\subsection{Démarche}

Le refactoring AoS → SoA consiste à remplacer le tableau d'objets \texttt{vector<ant>}
par trois tableaux distincts, un pour chaque attribut :

\begin{lstlisting}[language=C++, caption={Structure SoA dans ant\_colony (ant.hpp)}]
class ant_colony {
    std::vector<state_t>    m_states;    // etat de chaque fourmi
    std::vector<position_t> m_positions; // position (x, y)
    std::vector<std::size_t> m_seeds;   // graine aleatoire
public:
    void advance_all(pheronome& phen, const fractal_land& land,
                     const position_t& pos_food,
                     const position_t& pos_nest,
                     std::size_t& cpteur_food);
};
\end{lstlisting}

Cette organisation garantit que l'accès à un seul champ (par exemple les graines de tous
les agents) est \emph{contigu} en mémoire, ce qui favorise la prélecture matérielle
et la vectorisation par le compilateur.

\subsection{Résultats — version SoA (1 cœur)}

Les mesures suivantes proviennent de \texttt{src\_v2\_soa/} avec export CSV
et agrégation via \texttt{profile\_stats.exe}. Le run contient
\textbf{15847 itérations}.

Indicateurs métier observés :
\begin{itemize}
    \item première nourriture ramenée au nid à l'itération \textbf{2585} ;
    \item quantité finale de nourriture : \textbf{34312}.
\end{itemize}

\begin{table}[h]
\centering
\caption{Profiling par phase — version SoA, 1 cœur (run long)}
\label{tab:soa_timing}
\begin{tabular}{lrr}
\hline
\textbf{Phase}      & \textbf{Temps (ms/it)} & \textbf{Part (\%)} \\
\hline
Avancement fourmis  & 4.3421   & 3.67  \\
Évaporation         & 0.6928   & 0.58  \\
Update (swap)       & 0.0130   & 0.01  \\
Rendu SDL           & 113.4031 & 95.74 \\
\hline
\textbf{TOTAL}      & \textbf{118.4510} & 100.0 \\
\hline
\end{tabular}
\end{table}

\subsection{Discussion du surcoût SoA à 1 thread}

Comme pour AoS, le goulot dominant est le \textbf{rendu SDL} (\(\approx 95.7\%\)).
Le coût calculatoire SoA (hors rendu) est :
\[
T_{\mathrm{calc,SoA}} = 4.3421 + 0.6928 + 0.0130 = 5.0479\ \mathrm{ms/it}
\]
soit \(\approx 198\ \mathrm{it/s}\) en mode sans rendu.

Comparé au run AoS long (\(T_{\mathrm{calc,AoS}} = 4.8373\ \mathrm{ms/it}\)),
la version SoA est ici légèrement plus lente d'environ \(\frac{5.0479-4.8373}{4.8373}
\approx 4.35\%\). Cette observation reste cohérente avec un scénario où le bénéfice
de la réorganisation mémoire SoA se matérialise surtout en multi-thread, tandis que
le mono-cœur reste sensible aux surcoûts de structure et à la dynamique irrégulière
des déplacements.

En pratique, les conclusions robustes à ce stade sont :

\begin{enumerate}
    \item les mesures \emph{end-to-end} sont dominées par le rendu, donc peu discriminantes
        pour comparer finement AoS/SoA ;
    \item la comparaison AoS/SoA doit être faite prioritairement sur
        \(T_{\mathrm{calc}}\) (hors rendu) et sur les runs OpenMP multi-threads.
\end{enumerate}

Le gain SoA se matérialisera surtout en mode multi-thread, où la cohérence de cache
entre threads travaillant sur des fourchettes d'indices adjacentes est bien meilleure
qu'avec AoS.

% TODO : profil perf stat / valgrind cachegrind pour quantifier les défauts L1/L2
